%% Use the "normalphoto" option if you want a normal photo instead of cropped to a circle
% \documentclass[10pt,a4paper,normalphoto]{altacv}

\documentclass[10pt,a4paper,ragged2e,withhyper]{altacv}
%% AltaCV uses the fontawesome5 and packages.
%% See http://texdoc.net/pkg/fontawesome5 for full list of symbols.

% Change the page layout if you need to
\geometry{left=1.25cm,right=1.25cm,top=1.5cm,bottom=1.5cm,columnsep=1.2cm}

% The paracol package lets you typeset columns of text in parallel
\usepackage{paracol}

% Change the font if you want to, depending on whether
% you're using pdflatex or xelatex/lualatex
\ifxetexorluatex
  % If using xelatex or lualatex:
  \setmainfont{Roboto Slab}
  \setsansfont{Lato}
  \renewcommand{\familydefault}{\sfdefault}
\else
  % If using pdflatex:
  \usepackage[rm]{roboto}
  \usepackage[defaultsans]{lato}
  % \usepackage{sourcesanspro}
  \renewcommand{\familydefault}{\sfdefault}
\fi

% Change the colours if you want to
\definecolor{SlateGrey}{HTML}{2E2E2E}
\definecolor{LightGrey}{HTML}{666666}
\definecolor{DarkPastelRed}{HTML}{450808}
\definecolor{PastelRed}{HTML}{8F0D0D}
\definecolor{GoldenEarth}{HTML}{E7D192}
\colorlet{name}{black}
\colorlet{tagline}{PastelRed}
\colorlet{heading}{DarkPastelRed}
\colorlet{headingrule}{GoldenEarth}
\colorlet{subheading}{PastelRed}
\colorlet{accent}{PastelRed}
\colorlet{emphasis}{SlateGrey}
\colorlet{body}{LightGrey}

% Change some fonts, if necessary
\renewcommand{\namefont}{\Huge\rmfamily\bfseries}
\renewcommand{\personalinfofont}{\footnotesize}
\renewcommand{\cvsectionfont}{\LARGE\rmfamily\bfseries}
\renewcommand{\cvsubsectionfont}{\large\bfseries}


% Change the bullets for itemize and rating marker
% for \cvskill if you want to
\renewcommand{\itemmarker}{{\small\textbullet}}
\renewcommand{\ratingmarker}{\faCircle}

\begin{document}
\name{Tom De Caluwé}
\tagline{Embedded Software Engineer}
%% You can add multiple photos on the left or right
\photoR{3.1cm}{SAM_1688}
% \photoL{2.5cm}{Yacht_High,Suitcase_High}

\personalinfo{%
  % Not all of these are required!
  \email{decaluwe.t@gmail.com}
  \phone{+41 78 448 48 48}
  %\mailaddress{\r{A}smarkvegen 1206, 2365 \r{A}smarka, Norway}
  %\mailaddress{Liersesteenweg 99, 2640 Mortsel, Belgium}
  \mailaddress{Sonnenhofstrasse 7, 8853 Lachen SZ}
  %\location{Location, COUNTRY}
  %\homepage{www.homepage.com}
  %\twitter{@twitterhandle}
  \linkedin{tdecaluwe}
  \github{tdecaluwe}
  %\orcid{0000-0000-0000-0000}
  %% You can add your own arbitrary detail with
  %% \printinfo{symbol}{detail}[optional hyperlink prefix]
  % \printinfo{\faPaw}{Hey ho!}[https://example.com/]
  %% Or you can declare your own field with
  %% \NewInfoFiled{fieldname}{symbol}[optional hyperlink prefix] and use it:
  % \NewInfoField{gitlab}{\faGitlab}[https://gitlab.com/]
  % \gitlab{your_id}
  %%
  %% For services and platforms like Mastodon where there isn't a
  %% straightforward relation between the user ID/nickname and the hyperlink,
  %% you can use \printinfo directly e.g.
  % \printinfo{\faMastodon}{@username@instace}[https://instance.url/@username]
  %% But if you absolutely want to create new dedicated info fields for
  %% such platforms, then use \NewInfoField* with a star:
  % \NewInfoField*{mastodon}{\faMastodon}
  %% then you can use \mastodon, with TWO arguments where the 2nd argument is
  %% the full hyperlink.
  % \mastodon{@username@instance}{https://instance.url/@username}
}

\makecvheader
%% Depending on your tastes, you may want to make fonts of itemize environments slightly smaller
% \AtBeginEnvironment{itemize}{\small}

%% Set the left/right column width ratio to 6:4.
\columnratio{0.6}

% Start a 2-column paracol. Both the left and right columns will automatically
% break across pages if things get too long.
\begin{paracol}{2}

\cvsection{Experience}

\cvevent{Head of IT}{LeeGroup}{Sep 2024 -- Oct 2025}{Rothenburg, Switzerland}
\begin{itemize}
\item Building an IT department and infrastructure
\item Development projects in supply chain visibility, manufacturing resource planning and logistics optimization
\end{itemize}
\cvtag{python}
\cvtag{odoo}
\cvtag{shopify}

\divider

\cvevent{Python \& Javascript Developer}{Odoo}{Mar 2021 -- Feb 2024}{Ramilies, Belgium}
\begin{itemize}
\item Working on projects across a wide functional scope, including accounting, logistics, marketing automation and others
\item Contributing to the frontend framework for the odoo client
\end{itemize}
\cvtag{python}
\cvtag{javascript}
\cvtag{postgres}

\divider

\cvevent{IT Project Manager}{De Kampeerder}{May 2015 -- Nov 2019}{Antwerp, Belgium}
\begin{itemize}
\item Implementing and customizing Odoo as an ERP solution
\item Applying RFID technology in a retail environment to optimize logistical workflows
\item Integrating with online marketplaces
\end{itemize}
\cvtag{python}
\cvtag{java}
\cvtag{android}
\cvtag{rfid}

\divider

\cvevent{Software Developer}{Rubbens}{Project based (2016)}{Wichelen, Belgium}
\begin{itemize}
\item Developing a node.js webservice for \emph{Electronic Data Interchange} three-way match communications
\end{itemize}
\cvtag{node.js}
\cvtag{javascript}
\cvtag{edifact}

\divider

\cvevent{Dogsledding Guide}{Multiple Companies}{Dec 2019 - Apr 2024 (Seasonal)}{Canada, Norway \& Sweden}

% \divider

% \cvevent{Zipline Guide}{Zipzone Peachland}{Jun 2020 - Sep 2020}{Peachland, Canada}

\medskip

\cvsection{Volunteering}

\cvevent{Workshop Supervisor}{De Sleutel}{Sep 2014 - Mar 2015}{Antwerp, Belgium}

% use ONLY \newpage if you want to force a page break for
% ONLY the current column
%\newpage

%% Switch to the right column. This will now automatically move to the second
%% page if the content is too long.
\switchcolumn

\cvsection{Profile}

\begin{itemize}
\item Experience with both small and enterprise scale projects
\item Doesn't shy away from a challenge
\item Feels more at home in the backend, but knows his way around the frontend
\end{itemize}

\cvsection{Strengths}

\cvtag{problem solver}
\cvtag{eye for detail}\\
\smallskip
\cvtag{analytical mindset}

\cvsection{Technology}

\cvtag{python}
\cvtag{java}
\cvtag{javascript}
\cvtag{c++}\\
\smallskip
\cvtag{kotlin}
\cvtag{webgpu}

\divider

\cvtag{docker}
\cvtag{node.js}
\cvtag{android}
\cvtag{odoo}\\
\smallskip
\cvtag{git}
\cvtag{sql}
\cvtag{mongodb}
\cvtag{matlab}\\
\smallskip
\cvtag{sparql}
\cvtag{rfid}
\cvtag{arduino}

\cvsection{Languages}

\cvskill{English}{5}
\divider

\cvskill{French}{5}
\divider

\cvskill{Dutch}{5}
\divider

\cvskill{German}{2}

%% Yeah I didn't spend too much time making all the
%% spacing consistent... sorry. Use \smallskip, \medskip,
%% \bigskip, \vspace etc to make adjustments.
\medskip

\cvsection{Good to know}

\cvachievement{\faHeartbeat}{First Aid \& CPR}{Training \emph{Secourisme en mileu de travail} attended in french}

\divider

\cvachievement{\faCertificate}{Rock Climbing Instructor}{Certification course \emph{Initiator Rotsklimmen} organized by \emph{Sport Vlaanderen}}

\end{paracol}

\end{document}
